% Permite cortar las URL largas de la bibliografía también en los guiones
\PassOptionsToPackage{hyphens}{url}
\usepackage[utf8]{inputenc}
\usepackage[T1]{fontenc}
% english solo para partir bien las palabras del abstract; el idioma principal es spanish
\usepackage[english,spanish,es-tabla]{babel}
% El % sin espacio automático: el espacio se escribe a mano con \,
\spanishplainpercent
\decimalpoint

\usepackage{amsmath,amssymb,amsfonts}
\usepackage{graphicx}
\usepackage{geometry}
\usepackage{booktabs}
\usepackage{multirow}
\usepackage{array}
\usepackage{tabularx}
\usepackage{cite}
\usepackage{caption}
\usepackage{subcaption}
\usepackage{xcolor}
\usepackage{enumitem}
\usepackage{abstract}
\usepackage{titlesec}
\usepackage{verbatim} % Entorno comment que admite sangría antes de \end{comment}
\usepackage[colorlinks=true, linkcolor=blue, urlcolor=blue, citecolor=blue]{hyperref}

\graphicspath{{img/}}

% --- MÁRGENES ---
\geometry{
	letterpaper,
	top=2.0cm,
	bottom=2.0cm,
	left=1.8cm,
	right=1.8cm,
	columnsep=0.6cm
}

% --- FORMATO DE SECCIONES ---
\titleformat{\section}
{\normalfont\scshape\centering}
{\thesection.}
{0.5em}
{}

\renewcommand{\thesection}{\Roman{section}}

\titleformat{\subsection}
{\normalfont\itshape}
{\thesubsection.}
{0.5em}
{}

\renewcommand{\thesubsection}{\Alph{subsection}}

\titleformat{\subsubsection}[runin]
{\normalfont\itshape}
{}
{0em}
{}[.]

% --- CONFIGURACIÓN DEL RESUMEN ---
\renewcommand{\abstractname}{}
\setlength{\absleftindent}{0pt}
\setlength{\absrightindent}{0pt}

% Desactiva el título principal del resumen
\renewcommand{\abstractname}{}
\renewcommand{\absnamepos}{empty}
% Eliminar el título "Resumen" del entorno abstract por completo
\makeatletter
\renewenvironment{abstract}{%
	\small
	\par\noindent
}{%
	\par
}
\makeatother

% --- FIGURAS, TABLAS Y ALGORITMOS ---
\captionsetup{font=small, labelfont=bf, labelsep=period}
\captionsetup[sub]{font=footnotesize}
\setlist{nosep, leftmargin=*}

% Más espacio para flotantes en dos columnas
\renewcommand{\topfraction}{0.9}
\renewcommand{\bottomfraction}{0.8}
\renewcommand{\textfraction}{0.07}
\renewcommand{\floatpagefraction}{0.75}
\renewcommand{\dbltopfraction}{0.9}
\renewcommand{\dblfloatpagefraction}{0.75}
\setcounter{topnumber}{3}
\setcounter{dbltopnumber}{2}
\setcounter{totalnumber}{5}

\addto\captionsspanish{\renewcommand{\refname}{REFERENCIAS}}

% --- NOTACIÓN ---
% Arquitectura de la red: capas ocultas & capa de salida
\newcommand{\red}[2]{\mbox{#1}\,\&\,\mbox{#2}}
\newcommand{\onehot}{\textit{one-hot}}

% Paleta de los 18 estados de FLAME 3 (colores del simulador)
\definecolor{flame0}{HTML}{448CFD}\definecolor{flame1}{HTML}{17D9C7}\definecolor{flame2}{HTML}{2FF09A}
\definecolor{flame3}{HTML}{69FD65}\definecolor{flame4}{HTML}{A1FC3D}\definecolor{flame5}{HTML}{CDEB34}
\definecolor{flame6}{HTML}{EDCF39}\definecolor{flame7}{HTML}{FDA932}\definecolor{flame8}{HTML}{FA7A1F}
\definecolor{flame9}{HTML}{E84B0C}\definecolor{flame10}{HTML}{CB2B03}\definecolor{flame11}{HTML}{A01101}
\definecolor{flame12}{HTML}{000000}\definecolor{flame13}{HTML}{333333}\definecolor{flame14}{HTML}{666666}
\definecolor{flame15}{HTML}{999999}\definecolor{flame16}{HTML}{CCCCCC}\definecolor{flame17}{HTML}{FFFFFF}
\newcommand{\estadoflame}[2][black]{\fcolorbox{gray}{flame#2}{\makebox[0.9em]{\textcolor{#1}{\tiny #2}}}}
\newcommand{\leyendaflame}{%
	\par\smallskip
	{\scriptsize\setlength{\fboxsep}{0.8pt}%
	Temperatura: \estadoflame{0}\estadoflame{1}\estadoflame{2}\estadoflame{3}\estadoflame{4}\estadoflame{5}%
	\estadoflame{6}\estadoflame{7}\estadoflame{8}\estadoflame{9}\estadoflame[white]{10}\estadoflame[white]{11}%
	\quad Quemado: \estadoflame[white]{12}\estadoflame[white]{13}\estadoflame[white]{14}\estadoflame{15}%
	\estadoflame{16}\estadoflame{17}\par}}
% Versión en dos líneas para figuras de una columna
\newcommand{\leyendaflamecorta}{%
	\par\smallskip
	{\scriptsize\setlength{\fboxsep}{0.8pt}%
	Temperatura: \estadoflame{0}\estadoflame{1}\estadoflame{2}\estadoflame{3}\estadoflame{4}\estadoflame{5}%
	\estadoflame{6}\estadoflame{7}\estadoflame{8}\estadoflame{9}\estadoflame[white]{10}\estadoflame[white]{11}\\[1pt]
	Quemado: \estadoflame[white]{12}\estadoflame[white]{13}\estadoflame[white]{14}\estadoflame{15}%
	\estadoflame{16}\estadoflame{17}\par}}

% Paleta de los 7 estados de SEVIR
\definecolor{sevir0}{HTML}{F2F2F2}\definecolor{sevir1}{HTML}{9BE39B}\definecolor{sevir2}{HTML}{2E9E3A}
\definecolor{sevir3}{HTML}{F5D90A}\definecolor{sevir4}{HTML}{F28C1E}\definecolor{sevir5}{HTML}{D7191C}
\definecolor{sevir6}{HTML}{B10DC9}
\newcommand{\estadosevir}[2][black]{\fcolorbox{gray}{sevir#2}{\makebox[0.9em]{\textcolor{#1}{\tiny #2}}}}
\newcommand{\leyendasevir}{%
	\par\smallskip
	{\scriptsize\setlength{\fboxsep}{0.8pt}%
	Estado: \estadosevir{0}\estadosevir{1}\estadosevir[white]{2}\estadosevir{3}\estadosevir{4}%
	\estadosevir[white]{5}\estadosevir[white]{6}\par}}
